\section{Problema y motivacion}

% ---------------------------------------------------------------------------
% Marco a\~nadido el 2026-07-27. Faltaba lo m\'as importante: una respuesta directa a
% ``\textquestiondown{}aportamos algo?''. El deck contaba la historia del error durante
% 40 diapositivas y el aporte quedaba impl\'icito al final. Va aqu\'i, al principio, en
% positivo, y se cierra con lo mismo en la secci\'on 5.
% ---------------------------------------------------------------------------
\begin{frame}{\textquestiondown{}Qu\'e aportamos? --- la respuesta corta}
  \begin{block}{S\'i, y no es el art\'iculo que planeamos}
    \small
    No aportamos un descubrimiento de circuitos de regulaci\'on. \textbf{Aportamos una
    medici\'on de cu\'anto miente un protocolo de evaluaci\'on que este campo usa de
    forma rutinaria}, sobre un grafo biom\'edico real, con el protocolo corregido incluido.
  \end{block}

  \vspace{0.35em}
  \begin{columns}[T]
    \begin{column}{0.48\textwidth}
      \onslide<2->{%
      \begin{block}{Cuatro cosas concretas}
        \scriptsize
        \begin{enumerate}\setlength{\itemsep}{0.15em}
          \item \textbf{La inflaci\'on, cuantificada y descompuesta:} 0.99 $\to$ 0.63,
                y las dos causas se \textbf{potencian}.
          \item \textbf{Es del protocolo, no del modelo:} \textbf{6 arquitecturas} se
                inflan igual (0.24--0.31), \emph{una de ellas sin grafo}. El control sin
                entrenar no se mueve.
          \item \textbf{Un control que el campo no reporta:} 0/7 art\'iculos revisados
                publican una l\'inea base sin aprendizaje. Aqu\'i una heur\'istica de una
                l\'inea \textbf{le gana a todas}.
          \item \textbf{Un protocolo reutilizable,} con prueba autom\'atica de fuga.
        \end{enumerate}
      \end{block}}
    \end{column}
    \begin{column}{0.48\textwidth}
      \onslide<3->{%
      \begin{block}{Por qu\'e le importa a alguien m\'as}
        \scriptsize
        Los AUROC publicados en esta literatura viven entre \textbf{0.91 y 0.99}. Bajo ese
        mismo protocolo, una regla de una l\'inea que \textbf{ignora al miRNA} alcanza
        \textbf{0.87} --- dentro de esa banda.

        \vspace{0.25em}
        \emph{Un campo que nunca reporta un control sin modelo no puede saber si su 0.97 es
        un resultado o un efecto de popularidad.}
      \end{block}}

      \onslide<4->{%
      \begin{block}{Y lo que \emph{s\'i} funciona}
        \scriptsize
        Clasificaci\'on de tipo celular: \textbf{0.9916}.
      \end{block}}
    \end{column}
  \end{columns}
\end{frame}

\begin{frame}{Problema biom\'edico y computacional}
  \begin{block}{Contexto cl\'inico}
    La esclerosis m\'ultiple (EM) es una enfermedad autoinmune y neurodegenerativa cr\'onica con heterogeneidad celular, inflamatoria y temporal.
  \end{block}

  \vspace{0.5em}

  \begin{block}{Dificultad principal}
    Los miRNAs regulan genes y v\'ias de se\~nalizaci\'on de forma altamente dependiente del contexto celular. Los enfoques bulk promedian la se\~nal y ocultan qu\'e circuitos est\'an activos en monocitos, c\'elulas T, microgl\'ia u oligodendrocitos.
  \end{block}

  \vspace{0.5em}

  \begin{itemize}
    \item Un miRNA puede regular cientos de genes y un gen puede tener m\'ultiples reguladores.
    \item Las bases de interacci\'on por secuencia generan falsos positivos si no se incorpora contexto biol\'ogico.
    \item La EM requiere explicaciones por tipo celular, no solo m\'etricas globales.
  \end{itemize}
\end{frame}

\begin{frame}{Mensaje cl\'inico para abrir la defensa}
  \begin{block}{Problema que resuelve este trabajo}
    \small
    En EM sabemos que existen alteraciones transcript\'omicas e inmunes, pero no es trivial identificar qu\'e reguladores organizan esas alteraciones dentro de cada tipo celular.
  \end{block}

  \vspace{0.5em}

  \begin{columns}[T]
    \begin{column}{0.48\textwidth}
      \begin{block}{Aporte del proyecto}
        \small
        El proyecto integra c\'elulas, genes y miRNAs dentro de una sola estructura de grafo para convertir datos heterog\'eneos en hip\'otesis reguladoras priorizadas.
      \end{block}
    \end{column}
    \begin{column}{0.48\textwidth}
      \begin{block}{Valor para audiencia biom\'edica}
        \small
        El modelo no se limita a decir qu\'e c\'elula observa, sino que propone qu\'e miRNAs podr\'ian estar modulando genes clave en Th17, microgl\'ia, monocitos y oligodendrocitos.
      \end{block}
    \end{column}
  \end{columns}

  \vspace{0.4em}

  \begin{itemize}
    \item Frase sugerida para exponer: ``nuestro objetivo no es solo clasificar c\'elulas, sino priorizar circuitos reguladores con relevancia biol\'ogica en EM''.
  \end{itemize}
\end{frame}

\begin{frame}{Hip\'otesis y preguntas del proyecto}
  \begin{block}{Hip\'otesis}
    Integrar scRNA-seq, expresi\'on bulk de miRNAs y relaciones miRNA$\rightarrow$gen en un grafo heterog\'eneo permite aprender regulaci\'on espec\'ifica por tipo celular y priorizar circuitos con relevancia para EM.
  \end{block}

  \vspace{0.5em}

  \begin{enumerate}
    \item \textbf{Clasificaci\'on}: \textquestiondown{}Puede el modelo reconocer correctamente 11 tipos celulares?
    \item \textbf{Link prediction}: \textquestiondown{}Puede distinguir qu\'e pares miRNA-gen son plausibles en este contexto?
    \item \textbf{Interpretabilidad}: \textquestiondown{}Qu\'e miRNAs y circuitos emergen como candidatos biol\'ogicos prioritarios?
    \item \textbf{Traducci\'on}: \textquestiondown{}Qu\'e resultados son suficientemente estables para guiar validaci\'on externa?
  \end{enumerate}
\end{frame}

\begin{frame}{Objetivo general y entregables actuales}
  \begin{columns}[T]
    \begin{column}{0.57\textwidth}
      \begin{block}{Objetivo general}
        Construir una pipeline reproducible que una datos abiertos, genere un grafo heterog\'eneo y entrene una arquitectura tipo HGT para descubrir circuitos reguladores de miRNAs en EM.
      \end{block}

      \begin{block}{Entregables generados}
        \begin{itemize}
          \item Descarga y preprocesamiento de datos.
          \item Grafo PyG `HeteroData` listo para entrenamiento.
          \item Modelos baseline y variantes V2-V4.
          \item Interpretaci\'on: saliencia, circuitos top y enriquecimiento funcional.
          \item Figuras ya exportadas a PDF en `results/figures/`.
        \end{itemize}
      \end{block}
    \end{column}
    \begin{column}{0.40\textwidth}
      \begin{tikzpicture}[every node/.style={font=\small}]
        \node[draw=msblue, rounded corners, fill=mslight, text width=3.6cm, minimum height=0.95cm, align=center] (d1) at (0, 2.8) {scRNA-seq\\c\'elula \'unica};
        \node[draw=msblue, rounded corners, fill=mslight, text width=3.6cm, minimum height=0.95cm, align=center] (d2) at (0, 1.35) {miRNA bulk\\GEO / Affymetrix};
        \node[draw=msblue, rounded corners, fill=mslight, text width=3.6cm, minimum height=0.95cm, align=center] (d3) at (0, -0.1) {miRDB v6.0\\interacciones};
        \node[draw=msgreen, rounded corners, fill=green!8, text width=3.8cm, minimum height=1.1cm, align=center] (m) at (0, -1.95) {Grafo heterog\'eneo +\\miRNAGraphTransformer};
        \node[draw=msorange, rounded corners, fill=orange!10, text width=3.8cm, minimum height=1.1cm, align=center] (o) at (0, -3.8) {Tipos celulares +\\circuitos miRNA-gen};
        \draw[->, thick, msblue] (d1) -- (m);
        \draw[->, thick, msblue] (d2) -- (m);
        \draw[->, thick, msblue] (d3) -- (m);
        \draw[->, thick, msgreen] (m) -- (o);
      \end{tikzpicture}
    \end{column}
  \end{columns}
\end{frame}

\begin{frame}{Guion sugerido para una defensa de 10 minutos}
  \small
  \begin{tabularx}{\textwidth}{>{\bfseries}p{2.5cm}X}
    \toprule
    Tiempo & Punto a comunicar \\
    \midrule
    1 minuto & Plantear el problema cl\'inico y la necesidad de resoluci\'on por tipo celular \\
    2 minutos & Resumir datos y construcci\'on del grafo heterog\'eneo \\
    2 minutos & Explicar qu\'e aprende el modelo y por qu\'e no basta un enfoque cl\'asico \\
    2 minutos & Mostrar resultados V2 y comparativa contra baselines y ablaciones \\
    2 minutos & Presentar 2--3 hallazgos biol\'ogicos concretos \\
    1 minuto & Cerrar con limitaciones y siguiente validaci\'on \\
    \bottomrule
  \end{tabularx}
\end{frame}